\documentclass[10pt,a4paper]{article}
\usepackage{amsmath,amssymb,amsthm,graphicx}
\usepackage{hyperref}
\usepackage{enumitem}
\usepackage{mathtools}
\title{An exponential-time oracle blocks deterministic constructions while permitting quasipolynomial Kt approximation}
\author{Shuichi Hirahara, Zhenjian Lu, Hanlin Ren}
\date{doi:10.4230/LIPIcs.CCC.2023.6; CC BY 4.0}
\begin{document}
\maketitle
\noindent\textit{Editorial note.} Complete author published proof transcribed with original mathematical notation and exact theorem conclusions. The theorem statement prints omega(n), while the surrounding author prose/proof uses omega(1); both are retained as published, without silently editing the author. Display equations may be split across lines for page width; every author term is retained. Original last inequality in approximation proof ends in (1+epsilon)m and is retained.\par\medskip
\noindent\textit{Allowed standard prerequisites (editorial).}\begin{itemize}
\item Hirahara worst-case to average-case pseudorandom generator reconstruction. Original cited generator Lemma56 (cf. reference33) retained with complete parameters; standard earlier construction is not recursively reproved
\end{itemize}\medskip
\paragraph{Definitions and notations (original Sections 1.1, 2.1).}
For a complexity class $\mathfrak C$, a complexity-theoretic statement is $\mathfrak C$-relativizing if it is true relative to every oracle $\mathcal O\in\mathfrak C$.
\paragraph{Definition 20.} A language $L$ is in $\mathsf{EXPH}$ if there is a constant $k$, an exponential bound $e(n)=2^{\operatorname{poly}(n)}$, and a $\operatorname{poly}(e(n))$-time algorithm $M$ such that, for every input $x\in\{0,1\}^*$,
\[x\in L\Longleftrightarrow\exists y_1\forall y_2\exists y_3\forall y_4\ldots\exists y_k,\ M(x,y_1,\ldots,y_k)=1,\]
where $y_i\in\{0,1\}^{e(n)}$ for each $i\in[k]$. Equivalently, a language is in $\mathsf{EXPH}$ if it can be decided in exponential ($2^{\operatorname{poly}(n)}$) time with access to a $\mathsf{PH}$ oracle.
\paragraph{Pseudodeterministic algorithms.} A pseudodeterministic algorithm is a probabilistic algorithm that returns a fixed output with high probability on any given input. We will also consider pseudodeterministic oracle (or advice-taking) algorithms. Such an algorithm will maintain its pseudodeterministic behavior (i.e., outputting a fixed answer with high probability for any input) when a correct oracle (or advice) is given.
\paragraph{Dense property (original Section 1.3.2).} A property $Q\subseteq\{0,1\}^*$ is dense if for every input length $n\in\mathbb N$, we have that $|Q\cap\{0,1\}^n|\geq2^n/\operatorname{poly}(n)$.
\paragraph{Theorem 23 [68].} For every oracle $\mathcal O$, $\mathsf{BPP}^{\mathcal O}\subseteq(\Sigma^p_2)^{\mathcal O}$.

\paragraph{Time-Bounded Kolmogorov Complexity (original Section 2.1).}
We fix a time-optimal universal Turing machine $U$ and a time-optimal randomized Turing machine $V$.
\paragraph{Definition 21.} For a string $x\in\{0,1\}^*$, a time bound function $t:\mathbb N\to\mathbb N$, and an oracle $\mathcal O$, we define
\[K^{t,\mathcal O}(x):=\min_{d\in\{0,1\}^*}\{|d|:U^{\mathcal O}(d)\text{ outputs }x\text{ in at most }t(|x|)\text{ steps}\}.\]
\[rK^{t,\mathcal O}(x):=\min_{d\in\{0,1\}^*}\{|d|:V^{\mathcal O}(d)\text{ runs in at most }t(|x|)\text{ steps and }\Pr[V^{\mathcal O}(d)=x]\geq2/3\}.\]
We also consider Levin's notion of time-bounded Kolmogorov complexity.
\paragraph{Definition 22.} For a string $x\in\{0,1\}^*$ and an oracle $\mathcal O$, we define
\[\mathrm{Kt}^{\mathcal O}(x):=\min_{d\in\{0,1\}^*,t\in\mathbb N}\{|d|+\lceil\log t\rceil:U^{\mathcal O}(d)\text{ outputs }x\text{ in at most }t\text{ steps}\}.\]
\subsection*{5.2.2 Barriers for Deterministic Constructions and Lower Bounds for MKtP}
In this section, we construct an oracle in $\mathsf{EXP}$ relative to which there is no deterministic construction in $2^n/n^{\omega(1)}$ time (even infinitely often), showing that any non-trivial deterministic constructions using $\mathsf{PSPACE}$-relativizing techniques would separate $\mathsf{PSPACE}$ from $\mathsf{EXP}$.
As an interesting corollary, it is easy to approximate the Kt complexity to an $(1+\varepsilon)$ factor in this oracle world in deterministic $n^{O(\log n)}$ time. Therefore, although it is $\mathsf{EXP}$-hard to approximate the Kt complexity under $(\mathsf{P/poly})$-truth-table reductions and $\mathsf{NP}$-Turing reductions, any $\mathsf{PSPACE}$-relativizing proof that the Kt complexity requires deterministic $n^{\omega(\log n)}$ time to approximate would separate $\mathsf{PSPACE}$ from $\mathsf{EXP}$.
\paragraph{Definition 54.} For a constant $\varepsilon>0$, $\mathsf{Gap}_{\varepsilon}\mathsf{MKtP}$ is defined to be the promise problem $(\mathcal{YES}_n,\mathcal{NO}_n)_{n\in\mathbb N}$, where
\begin{align*}
\mathcal{YES}_n&:=\{(x,s)\in\{0,1\}^n\times\mathbb N:\mathrm{Kt}(x)\leq s\},\\
\mathcal{NO}_n&:=\{(x,s)\in\{0,1\}^n\times\mathbb N:\mathrm{Kt}(x)>(1+\varepsilon)\cdot s\}.
\end{align*}
Our proof relies heavily on the equivalence between non-trivial derandomization and the hardness of $\mathsf{Gap}_{\varepsilon}\mathsf{MKtP}$ [34]. It is not hard to construct an oracle under which non-trivial derandomization is impossible, which means that a dense subset of the complement of $\mathsf{MKtP}$ can be accepted by an efficient algorithm $A$. Then, we use the worst-case to average-case reduction of [33] to transform $A$ into a worst-case approximation algorithm for $\mathsf{Gap}_{\varepsilon}\mathsf{MKtP}$.

\paragraph{Theorem 55.} There exists an oracle $\mathcal O\in\mathsf{EXP}$ such that
\begin{enumerate}
\item there is a dense property $Q\in\mathsf P^{\mathcal O}$ such that every deterministic algorithm that runs in time $2^n/n^{\omega(n)}$ fails to find a string in $Q\cap\{0,1\}^n$ on almost every input length $n$, and
\item $\mathsf{Gap}_{\varepsilon}\mathsf{MKtP}^{\mathcal O}\in\mathsf{DTIME}^{\mathcal O}[n^{O(\log n)}]$ for every constant $\varepsilon>0$.
\end{enumerate}

To show the worst-case to average-case reduction, we use the following pseudorandom generator construction.
\paragraph{Lemma 56 (cf. [33]).} For any $d,m\leq2^n$, $\varepsilon>0$, there exists a ``pseudorandom generator construction''
\[G:\{0,1\}^n\times\{0,1\}^d\to\{0,1\}^m\]
such that for any distinguisher $D:\{0,1\}^m\to\{0,1\}$, if
\[\Pr_{w\sim\{0,1\}^m}[D(w)=1]-\Pr_{z\sim\{0,1\}^d}[D(G(x,z))=1]\geq\frac12,\]
then
\[K^{\operatorname{poly}(n),D}(x)\leq\exp(\ell^2/d)\cdot m+d+O(\log n),\]
where $\ell=O(\log n)$. Moreover, $G(x,z)$ can be computed in time $\operatorname{poly}(n)$.

\begin{proof}[Proof of Theorem 55]
We construct an oracle $\mathcal O\in\mathsf{EXP}$ under which a dense subset of $\mathsf{MKtP}$ is in $\mathsf P$. We start with $n:=1$, $\mathcal O:=\emptyset$, and $F:=\emptyset$, where $F$ is a set of ``frozen'' strings. The construction of $\mathcal O$ in Stage $n$ is as follows. Let $\mathcal O_n$ denote the state of the oracle $\mathcal O$ at the beginning of stage $n$. Let $E_n$ be the set of all strings $x\in\{0,1\}^n$ such that $\mathrm{Kt}^{\mathcal O_n}(x)\leq n-c\log n$, where $c$ is a sufficiently large constant (e.g., $c:=3$). Let $F_n$ be the set of all strings $q\in\{0,1\}^*$ such that there exist $k\in[n-c\log n]$ and a description $d\in\{0,1\}^k$ of a Turing machine such that the universal Turing machine $U^{\mathcal O_n}$ on input $d$ makes the query $q$ in time $2^{n-k-c\log n}$. Then, we update $F:=F\cup F_n$ and $\mathcal O_{n+1}:=\mathcal O_n\cup(\{0,1\}^n\setminus(E_n\cup F))$ and move on to the next stage $n+1$. This completes the description of the oracle $\mathcal O:=\bigcup_{n\in\mathbb N}\mathcal O_n$. It is easy to observe that $\mathcal O\in\mathsf{EXP}$.

The oracle $\mathcal O$ is a dense subset of $\{x:\mathrm{Kt}^{\mathcal O}(x)>n-c\log n\}$ in the following sense:
\begin{enumerate}
\item $x\notin\mathcal O$ for every string $x\in\{0,1\}^n$ such that $\mathrm{Kt}^{\mathcal O}(x)\leq n-c\log n$.
\item $x\in\mathcal O$ for at least half of the strings $x\in\{0,1\}^n$.
\end{enumerate}
To see the first property, it suffices to show that if $x\in\{0,1\}^n$ satisfies $\mathrm{Kt}^{\mathcal O}(x)\leq n-c\log n$, then $\mathrm{Kt}^{\mathcal O_n}(x)\leq n-c\log n$ (which implies $x\in E_n$ and thus $x\notin\mathcal O$). Assuming that there exists a string $d$ such that $U^{\mathcal O}(d)$ outputs $x$ in time $2^{n-c\log n-|d|}$, we claim that $U^{\mathcal O_n}(d)$ also outputs $x$. If not, there exists a query in $\mathcal O\setminus\mathcal O_n$ made by $U^{\mathcal O}(d)$. Let $q$ be the first query in $\mathcal O\setminus\mathcal O_n$. The same query is also queried during the computation of $U^{\mathcal O_n}(d)$. Thus, $q\in F_n$, which implies $q\notin\mathcal O$, which is a contradiction.

To see the second property, we bound the size of $E_n\cup F$ at Stage $n$. The size of $E_n$ is at most $\sum_{k\in[n-c\log n]}2^k\leq2^n/n^{c-1}$. The number $|F_n|$ of frozen strings in Stage $n$ is at most
\[\sum_{k\in[n-c\log n]}2^k\cdot2^{n-k-c\log n}\leq2^n/n^{c-1}.\]
Thus, the number $|F|$ of frozen strings until Stage $n$ is at most $\sum_{n'=1}^n|F_{n'}|\leq2^n/n^{c-2}$. Overall, we obtain
\[|\mathcal O\cap\{0,1\}^n|\geq2^n-|E_n|-|F|\geq\tfrac12\cdot2^n.\]
Let $Q:=\mathcal O$. We prove that $Q\in\mathsf P^{\mathcal O}$ is a dense property such that every $2^n/n^{\omega(1)}$-time deterministic algorithm $A$ fails to find a string in $Q\cap\{0,1\}^n$ on input $1^n$. By the definition of Kt, we have $\mathrm{Kt}^{\mathcal O}(A(1^n))\leq O(\log n)+\log(2^n/n^{\omega(1)})\leq n-\omega(\log n)$. Thus, the output $A(1^n)$ of $A$ is not in $\mathcal O$.

We now use the worst-case to average-case reduction of [34] to obtain an approximation algorithm $A^{\mathcal O}$ for $\mathsf{Gap}_{\varepsilon}\mathsf{MKtP}$. Let $x\in\{0,1\}^n$ and $s\in\mathbb N$ be an instance of $\mathsf{Gap}_{\varepsilon}\mathsf{MKtP}$. Let $d=k\cdot O(\log^2 n)$, where $k=k(\varepsilon)$ is a sufficiently large constant that will be chosen depending on $\varepsilon>0$. We may assume without loss of generality that $O(k^2\cdot\log^2 n)\leq s\leq n+O(\log n)$ because whether $\mathrm{Kt}(x)\leq O(k^2\cdot\log^2 n)$ or not can be decided in time $n^{O(k^2\cdot\log n)}$ by an exhaustive search. Let $m$ be a parameter chosen later. Let $G$ be the pseudorandom generator construction of Lemma 56. The algorithm $A^{\mathcal O}$ accepts $(x,s)$ if and only if $G(x,z)\notin\mathcal O$ for every $z\in\{0,1\}^d$. The running time of $A^{\mathcal O}$ is $2^d\operatorname{poly}(n)=n^{O(\log n)}$.

We prove the correctness of $A^{\mathcal O}$. Assume that $\mathrm{Kt}^{\mathcal O}(x)\leq s$. Then, for every $z\in\{0,1\}^d$, we have
\[\mathrm{Kt}^{\mathcal O}(G(x,z))\leq\mathrm{Kt}^{\mathcal O}(x)+O(d+\log n)\leq s+O(d+\log n)\leq m-c\log m,\]
where the last inequality holds by choosing a sufficiently large $m=s+O(d+\log n)$. It follows that $G(x,z)\notin\mathcal O$ and that $A^{\mathcal O}$ accepts $(x,s)$. Conversely, assume that $A^{\mathcal O}$ accepts $(x,s)$, which means that $\Pr_{z\sim\{0,1\}^d}[\mathcal O(G(x,z))=1]=0$. We claim that $\mathrm{Kt}^{\mathcal O}(x)\leq(1+\varepsilon)\cdot s$. Since $\mathcal O$ is dense, we have
\[\Pr_{w\sim\{0,1\}^m}[\mathcal O(w)=1]-\Pr_{z\sim\{0,1\}^d}[\mathcal O(G(x,z))=1]\geq\tfrac12.\]
By Lemma 56, we obtain
\begin{multline*}
K^{\operatorname{poly}(n),\mathcal O}(x)\leq\exp(O(\log^2 n)/d)\cdot m+d+O(\log n)\\
\leq\exp(1/k)\cdot m+O(k\cdot\log^2 n).
\end{multline*}
In particular,
\begin{multline*}
\mathrm{Kt}^{\mathcal O}(x)\leq\exp(1/k)\cdot m+O(k\cdot\log^2 n)\\
\leq(1+O(1/k))\cdot m+O(k\cdot\log^2 n)\leq(1+\varepsilon)\cdot m,
\end{multline*}
where we choose a sufficiently large $k:=O(1/\varepsilon)$ and use that $m\geq O(k^2\cdot\log^2 n)$ in the last inequality.
\end{proof}

\begin{thebibliography}{99}
\bibitem[31]{GIPS} Shafi Goldwasser, Russell Impagliazzo, Toniann Pitassi, and Rahul Santhanam. On the pseudo-deterministic query complexity of NP search problems. CCC2021, LIPIcs200,36:1--36:22. doi:10.4230/LIPIcs.CCC.2021.36.
\bibitem[33]{H18} Shuichi Hirahara. Non-black-box worst-case to average-case reductions within NP. FOCS2018,247--258. doi:10.1109/FOCS.2018.00032.
\bibitem[34]{H20} Shuichi Hirahara. Non-disjoint promise problems from meta-computational view of pseudorandom generator constructions. CCC2020,LIPIcs169,20:1--20:47. doi:10.4230/LIPIcs.CCC.2020.20.
\bibitem[68]{Sipser} Michael Sipser. A complexity theoretic approach to randomness. STOC1983,330--335. doi:10.1145/800061.808762.
\end{thebibliography}
\end{document}
